% Generated by t3_write_results.m; numerical values are not edited manually.
The example uses three tokens with three features each, one attention head, and a shared $3\to6\to3$ feed-forward network. All projection and feed-forward parameters are hand-chosen constants; no training is performed. Attention is bidirectional, with no mask or positional encoding. The three levels are functional stages within one head, not three stacked Transformer blocks.

\begin{table}[htbp]\centering\small
\renewcommand{\arraystretch}{1.35}
\begin{tabular}{lc}\toprule
Quantity & Numerical value \\
\midrule
Input $X$ & $\begin{bmatrix}1.0000 & 0.2000 & -0.4000 \\ -0.3000 & 0.8000 & 0.5000 \\ 0.6000 & -0.5000 & 0.9000\end{bmatrix}$ \\[0.5em]
Attention program $A$ & $\begin{bmatrix}0.3776 & 0.2084 & 0.4141 \\ 0.3928 & 0.3718 & 0.2354 \\ 0.2188 & 0.4007 & 0.3805\end{bmatrix}$ \\[0.5em]
Head output $Z$ & $\begin{bmatrix}0.6542 & -0.0256 & 0.2184 \\ 0.4430 & 0.1450 & 0.2536 \\ 0.4287 & -0.0404 & 0.4011\end{bmatrix}$ \\[0.5em]
Block output $Y$ & $\begin{bmatrix}1.3174 & -0.4559 & -0.9535 \\ -1.3294 & 0.9408 & 0.5024 \\ 0.8290 & -1.4501 & 0.5075\end{bmatrix}$ \\[0.5em]
\bottomrule\end{tabular}
\caption[Three-token example: input, attention, head, and block outputs]{Numerical matrices for the three-token example. The fixed projections are defined in Equation~\eqref{eq:t3-projections}; $A$ and $Z$ follow Equations~\eqref{eq:t3-attention} and~\eqref{eq:t3-execution}. $Y$ includes the output projection, residual paths, row-wise LayerNorm, and shared feed-forward sublayer (Equations~\eqref{eq:t3-postnorm} through~\eqref{eq:t3-block-output}). Values are rounded to four decimals here; full precision is retained in the CSV/JSON/MAT exports. MATLAB producers: \path{run_t3_example.m} (computation) and \path{t3_write_results.m} (this table).}
\label{tab:t3-numeric}\end{table}

A separate illustrative vocabulary head maps each three-feature block output to four symbol probabilities. Its weights are hand-chosen and untrained; these probabilities do not make the example a trained language model.

\begin{table}[htbp]\centering\small
\begin{tabular}{lrrrrl}\toprule
Token & alpha & beta & gamma & delta & Prediction \\
\midrule
$t_1$ & 0.3671 & 0.0626 & 0.4338 & 0.1365 & gamma \\
$t_2$ & 0.0762 & 0.6240 & 0.1586 & 0.1412 & beta \\
$t_3$ & 0.4304 & 0.0300 & 0.0916 & 0.4479 & delta \\
\bottomrule\end{tabular}
\caption[Illustrative four-symbol vocabulary readout]{Final symbol probabilities from the linear $3\to4$ readout and row-wise softmax in Equation~\eqref{eq:t3-readout}. Columns index vocabulary symbols, whereas columns of the attention matrix index source tokens. This head is illustrative and untrained. MATLAB producers: \path{t3_vocabulary_readout.m} (computation) and \path{t3_write_results.m} (this table).}\label{tab:t3-readout}\end{table}

The independent scalar-loop and vectorized implementations agree to a maximum absolute difference of $1.33\times 10^{-15}$. The checks also verify positive normalized attention rows, query and context effects with fixed parameters, controlled key/attention program swaps, token-permutation equivariance, feature-wise LayerNorm, and token-wise feed-forward locality.
This example uses exact scalar multiplication. It does not implement or validate the historical approximate CTRNN multiplier, and it makes no active-inference claim.
